% !TEX TS-program = lualatex

%%%
% src::
%     url = https://github.com/pfradin/luadraw/discussions/89#discussioncomment-14457144
%%%

\documentclass[border = 3pt]{standalone}

\usepackage[svgnames]{xcolor}
\usepackage[3d]{luadraw}

\begin{document}

\begin{luadraw*}{name = param-surf-colmap-OK}
local ld = luadraw

local M      = ld.pt3d.M
local Origin = ld.pt3d.Origin
local vecJ   = ld.pt3d.vecJ
local vecK   = ld.pt3d.vecK

local palext = require 'luadraw_palettes'

local pal = palext.Gnuplot

local pi, sin, cos, tanh = math.pi, math.sin, math.cos, math.tanh

------
-- Nous gardons juste un quart de la ¨surf, un morceau n'ayant
-- que des faces orientées vers l'observateur.
------
local projplane = ld.surface(
  function(u, v)
    return M(
      2*(1 + cos(v))*cos(u),
      2*(1 + cos(v))*sin(u),
      2*sin(v)*tanh(u-pi)
    )
  end,
  pi, 0, 0, pi,
  {25, 15}
)

------
-- Le reste de la surface s'obtient via des symétries, mais en
-- prenant garde à donner la bonne orientation aux faces visibles.
------
projplane = ld.concat(
  projplane,
  ld.reverse_face_orientation(
    ld.sym3d(
      projplane,
      {Origin, vecJ}
    )
  )
)

projplane = ld.concat(
  projplane,
  ld.reverse_face_orientation(
    ld.sym3d(
      projplane,
      {Origin, vecK}
    )
  )
)

local graphview = ld.graph3d:new{
  viewdir = {120, 55},
  bbox    = false
}

graphview:Dpoly(
  ld.facet2poly(projplane),
  {
    usepalette = {pal, "z"},
    mode       = ld.mShaded
  }
)

graphview:Show()
\end{luadraw*}

\end{document}
